\chapter*{Preface}
\addcontentsline{toc}{chapter}{Preface}
\markboth{Preface}{Preface}

\section*{Who this manual is for}

This manual is written for the person who installs, configures and
runs a mybibli instance for their household, a small association,
or a community library. It assumes you can:

\begin{itemize}
  \item open a terminal and run \texttt{docker compose} commands;
  \item edit a plain-text configuration file with the editor of
        your choice;
  \item read a URL printed in a log line and open it in a browser.
\end{itemize}

No Rust, MariaDB or HTMX knowledge is required. The manual is
deliberately scoped to the \emph{operator} role — building mybibli
from source or modifying its code is covered by
\texttt{CLAUDE.md} and the planning artifacts in the repository.

\section*{Conventions}

\textbf{Commands.} Shell snippets are shown in a fixed-width box:

\begin{lstlisting}
docker compose up -d
\end{lstlisting}

\textbf{Configuration files.} File excerpts use the same style with
the file name called out in the surrounding text. Comments inside
the snippet use the file's native comment marker.

\textbf{Callouts.} Three sidebars highlight important text:

\begin{note}
A \emph{note} provides background or context that helps you understand
why something works the way it does. Skipping it is safe.
\end{note}

\begin{tip}
A \emph{tip} suggests a practice that makes day-to-day use smoother.
Try it once you are comfortable with the basics.
\end{tip}

\begin{warning}
A \emph{warning} flags a step where a mistake can lose data, lock
you out of the admin panel, or expose the instance to unauthorized
access. Read these in full.
\end{warning}

\textbf{Cross-references.} Chapters and sections are linked in the
PDF — click any blue title or page number to jump there. The
\hyperref[index]{index} at the end maps every important term to its
page of first use.

\section*{Getting help}

The canonical source of truth for mybibli is the GitHub repository:

\begin{center}
\url{https://github.com/guycorbaz/mybibli}
\end{center}

Open an issue against the \texttt{guycorbaz/mybibli} repository when
you hit a bug, a documentation gap, or a missing feature. The issue
templates ask for the version (visible at the bottom of any page in
the admin panel), your install method, and the steps to reproduce.

\section*{Manual coverage}

This manual covers mybibli \textbf{1.1.0 and later}. Older versions
shipped a security-relevant bug where the first-launch setup wizard
was silently bypassed — see chapter~\ref{ch:upgrade} for the
upgrade path from 1.0.x.

\vspace{1cm}
\hrule
\vspace{0.5cm}
\textit{Happy cataloging.}
